% Chapter 13: Team Patterns & Governance
\chapter{Team Patterns and Governance}
\label{ch:team}

\begin{itshape}
You have been using Claude effectively for months.
Your team lead asks: ``Can we standardize this across our 8-person team?''
The answer is yes---but individual mastery does not automatically
translate to team productivity. Shared conventions require infrastructure.
\end{itshape}

\section{Shared CLAUDE.md in Version Control}
\label{sec:shared-claude-md}

\marginnote[Official]{Claude Code included with Team plan seats.
Shared CLAUDE.md in VCS.\sourceurl{https://support.claude.com/en/articles/11845131-use-claude-code-with-your-team-or-enterprise-plan}}

The project-level \code{CLAUDE.md} (structured per \cref{ch:architecture})
and \code{.claude/\allowbreak{}settings.json} are committed to version control.
Every team member gets the same project context, the same permissions,
and the same hooks.

\practitioner{Three files form the team standard:}

\begin{description}
  \item[\code{CLAUDE.md}] Project knowledge: build commands, architecture, conventions.
    Reviewed in PRs like any other code.
  \item[\code{.claude/settings.json}] Team permissions: deny rules, required hooks.
    Committed to git.
  \item[\code{.claude/settings.local.json}] Personal preferences: model choice,
    output style, personal API tokens. Auto-gitignored.
\end{description}

\begin{whybox}[Why Separate Committed and Local Settings]
If model choice is in \code{settings.json}, developers who need Opus
for complex tasks must override team settings.
If it is in \code{settings.local.json}, each developer chooses their own
model without affecting others.

The principle: team standards (deny rules, hooks) are committed.
Personal preferences (model, style) are local.
\end{whybox}

% -------------------------------------------------------------------
\section{Managed Settings for IT}
\label{sec:managed-settings}

\official{IT administrators can deploy \code{managed-settings.json}
files that override all other settings.} This enables:

\begin{itemize}
  \item Organization-wide deny rules (sensitive file access)
  \item Mandatory hooks (security scanning, compliance checks)
  \item Approved MCP server lists
  \item Model restrictions (e.g., Sonnet only for cost control)
\end{itemize}

Managed settings deploy to system-level paths and cannot be overridden
by project or user settings (see \cref{sec:settings} for the full 5-level
precedence hierarchy). They are the IT equivalent of hooks---non-negotiable.

% -------------------------------------------------------------------
\section{Onboarding New Team Members}
\label{sec:onboarding}

\practitioner{Day 1 with Claude for a new team member:}

\begin{enumerate}
  \item Clone the repo. CLAUDE.md and settings load automatically.
  \item Run \code{/init} to verify Claude understands the project.
    Compare output to the existing CLAUDE.md---they should align.
  \item Create \code{.claude/settings.local.json} for personal preferences.
  \item Run the team's \code{/check} command to verify the environment works.
  \item Review the team's hook configuration:
    what runs on file write? On commit? Why?
\end{enumerate}

The goal: a new developer should be productive with Claude
within their first day, without reading a setup guide.
The configuration in version control should be self-documenting.

% -------------------------------------------------------------------
\section{Code Review Patterns with Claude}
\label{sec:code-review}

\practitioner{Two patterns for team code review with Claude
(see also the Writer/Reviewer pattern in \cref{sec:writer-reviewer}):}

\subsection{Pre-PR Self-Review}

Before creating a PR, have Claude review your own changes:

\begin{minted}{text}
Review the diff in `git diff main...HEAD` for:
1. Logic errors and bugs
2. Missing test coverage
3. Style violations per our CLAUDE.md
4. Security concerns
Report issues with file:line references and severity.
\end{minted}

This catches obvious issues before reviewers see the code,
reducing review cycles.

\subsection{Reviewer-Assist}

During review, use Claude to investigate specific concerns:

\begin{minted}{text}
In PR #234, the author changed the feature
normalization strategy in
src/features/preprocessing.py.
What are the implications for models trained on
the old normalization? Check if any downstream
models depend on the old feature distributions.
\end{minted}

\margintip{The reviewer still makes the decision.
Claude provides the analysis.
AI augments review; it does not replace it.}

% -------------------------------------------------------------------
\section{Quick Reference}

\begin{itemize}
  \item Commit \code{CLAUDE.md} and \code{settings.json} for team standards
  \item Keep \code{settings.local.json} for personal preferences (auto-gitignored)
  \item Managed settings for IT policy enforcement (highest precedence)
  \item New team members productive on day 1 via committed configuration
  \item Pre-PR self-review catches issues before human reviewers see them
\end{itemize}

% -------------------------------------------------------------------
\begin{trybox}
Draft a team \code{settings.json} with:
\begin{enumerate}
  \item 3 deny rules protecting sensitive files (\code{.env*}, \code{secrets/}, credentials)
  \item 1 ask rule for destructive operations (\code{git push})
  \item 1 PostFileWrite hook for auto-formatting
  \item 1 PreCommit hook for your test suite
\end{enumerate}
Have a teammate review it. Does it match their expectations?
\end{trybox}
